% =============================================================================
% CAPA COMPLETA (BROCHURA / PAPERBACK) — MOLAMBUDOS · EDIÇÃO PORTUGUÊS (PT-BR) · 160 x 230 mm
% Gerada no R399. Toda a geometria é DERIVADA dos parâmetros abaixo: para
% ajustar, mude só os \def e recompile — nenhum número está fixo no corpo.
%
% Por que um arquivo novo em vez de editar capa_completa.tex: aquele é o
% projeto de CAPA DURA (case laminate, wrap de 0.591in, painéis 6.197x9.236in
% ≈ 157x235mm, correspondendo ao trim KDP hardcover 6.14x9.21in). A KDP não
% oferece 160x230mm em capa dura; em brochura o trim personalizado é aceito.
% Os dois projetos coexistem — nada foi perdido.
%
% CONFERIR ANTES DE ENVIAR À GRÁFICA:
%   1. \Paginas tem de ser a contagem final do miolo aprovado.
%   2. \EspessuraPapel é a constante da KDP por tipo de impressão:
%        0.002252 in/pág — P&B em papel branco
%        0.0025   in/pág — P&B em papel creme
%        0.002347 in/pág — cor (premium) em papel branco   <-- em uso
%      O miolo tem fundo sépia em todas as páginas, logo é impressão em cor.
%   3. Baixar o template oficial da KDP para o trim e a paginação finais e
%      conferir a lombada contra ele. A KDP recusa capa cuja lombada não
%      corresponda ao número de páginas.
% =============================================================================

\documentclass[12pt]{article}

% --- PARÂMETROS -------------------------------------------------------------
\def\TrimLargura{6.29921}   % 160 mm em polegadas
\def\TrimAltura{9.05512}    % 230 mm em polegadas
\def\Paginas{435}           % miolo main_kdp_pt_160x230mm.pdf
\def\EspessuraPapel{0.002347}
\def\Sangria{0.125}         % bleed KDP para brochura
\def\MargemSegura{0.25}     % recuo mínimo do conteúdo em relação ao corte

\usepackage{tikz}
\usetikzlibrary{calc}
\usepackage{pgf}

% --- GEOMETRIA DERIVADA -----------------------------------------------------
\pgfmathsetmacro{\Lombada}{\Paginas*\EspessuraPapel}
\pgfmathsetmacro{\CapaLargura}{2*\Sangria + 2*\TrimLargura + \Lombada}
\pgfmathsetmacro{\CapaAltura}{2*\Sangria + \TrimAltura}
\pgfmathsetmacro{\ContraX}{\Sangria}
\pgfmathsetmacro{\LombadaX}{\Sangria + \TrimLargura}
\pgfmathsetmacro{\FrenteX}{\Sangria + \TrimLargura + \Lombada}
\pgfmathsetmacro{\CorteTopo}{\Sangria + \TrimAltura}
\pgfmathsetmacro{\PainelLargura}{\TrimLargura + \Sangria}
\pgfmathsetmacro{\LombadaCentro}{\Sangria + \TrimLargura + \Lombada/2}
% A reserva do código de barras fica no canto inferior esquerdo da quarta capa
% e é dimensionada para COBRIR o código de barras de teste que está embutido em
% contracapa.png (ISBN fictício 978-65-01-23456-7, que não corresponde ao ISBN
% real desta obra). O ideal é remover aquele elemento da arte de origem; até lá,
% esta reserva o oculta e cumpre a exigência da KDP de deixar a área livre.
\pgfmathsetmacro{\CodigoX}{\Sangria + 0.10}
\pgfmathsetmacro{\CodigoY}{\Sangria + 0.02}
\def\CodigoLargura{2.15in}
\def\CodigoAltura{0.95in}
\pgfmathsetmacro{\CorteDireita}{\CapaLargura - \Sangria}
\pgfmathsetmacro{\MeiaAltura}{\CapaAltura/2}
\pgfmathsetmacro{\LogoY}{\Sangria + 0.22}

\usepackage[
  paperwidth=\CapaLargura in,
  paperheight=\CapaAltura in,
  margin=0in
]{geometry}

\usepackage[T1]{fontenc}
\usepackage[utf8]{inputenc}
\usepackage[brazilian]{babel}
\usepackage{ebgaramond}
\usepackage{graphicx}
\usepackage{xcolor}

\pagestyle{empty}

% --- Identidade visual (idêntica ao projeto de capa dura) --------------------
\definecolor{fundoescuro}{HTML}{1A1410}
\definecolor{sepiabg}{HTML}{F2E8CF}
\definecolor{customgold}{HTML}{CBA135}
\definecolor{sangue}{HTML}{8B0000}
\definecolor{textwhite}{HTML}{F5F0EB}

\begin{document}
\begin{tikzpicture}[remember picture, overlay]
% Tudo é desenhado dentro de um scope ancorado no canto inferior esquerdo da
% PÁGINA. Sem esse shift, `overlay` faz as coordenadas partirem do ponto de
% texto corrente — foi exatamente o que deixou capa_completa.pdf (o projeto de
% capa dura) com todo o conteúdo empurrado para fora da página.
\begin{scope}[shift={(current page.south west)}, x=1in, y=1in]

  % 1. Fundo integral, sangrado
  \fill[fundoescuro] (0,0) rectangle (\CapaLargura, \CapaAltura);

  % 2. Quarta capa — a arte cobre o painel inteiro mais a sangria externa
  \begin{scope}
    \clip (0, 0) rectangle (\LombadaX, \CapaAltura);
    \node[anchor=south west, inner sep=0pt] at (0, 0) {%
      \includegraphics[width=\PainelLargura in, height=\CapaAltura in]{frontmatter/contracapa.png}%
    };
  \end{scope}

  % 2b. Reserva do código de barras (KDP: 2.0 x 1.2 in, livre e clara)
  \node[anchor=south west, fill=white, inner sep=0pt, minimum width=\CodigoLargura,
        minimum height=\CodigoAltura] at (\CodigoX, \CodigoY) {%
    \begin{tabular}{c}
      {\color{black}\sffamily\bfseries\small ISBN 9798189170492}\\[0.5em]
      {\color{black!40}\sffamily\tiny área reservada ao código de barras}
    \end{tabular}%
  };

  % 3. Lombada
  \fill[fundoescuro!95!black] (\LombadaX, 0) rectangle (\FrenteX, \CapaAltura);

  \node[anchor=center, rotate=-90] at (\LombadaCentro, \MeiaAltura) {%
    \begin{minipage}{7.6in}
      \centering
      {\color{customgold}\fontfamily{ebgaramond}\selectfont\bfseries\Large MOLAMBUDOS}
      \hspace{0.8em}{\color{textwhite}\small\textbullet}\hspace{0.8em}
      {\color{textwhite}\fontfamily{ebgaramond}\selectfont\itshape\small O Diário do Paciente 1.260}
      \hspace{1.5em}
      {\color{customgold!80}\fontfamily{ebgaramond}\selectfont\scriptsize Marcelo Dias de Carvalho Filho \& Marcelo Claro}
    \end{minipage}%
  };

  \node[anchor=south] at (\LombadaCentro, \LogoY) {%
    \includegraphics[width=0.42in]{frontmatter/logo-white.png}%
  };

  % 4. Capa frontal — cobre o painel mais a sangria externa
  \begin{scope}
    \clip (\FrenteX, 0) rectangle (\CapaLargura, \CapaAltura);
    \node[anchor=south west, inner sep=0pt] at (\FrenteX, 0) {%
      \includegraphics[width=\PainelLargura in, height=\CapaAltura in]{frontmatter/capa.png}%
    };
  \end{scope}

  % 5. Guias técnicas (corte e dobras da lombada) — remover antes do envio
  \draw[white!25, dotted, line width=0.3pt]
    (\Sangria, \Sangria) rectangle (\CorteDireita, \CorteTopo);
  \draw[customgold!45, dotted, line width=0.5pt] (\LombadaX, 0) -- (\LombadaX, \CapaAltura);
  \draw[customgold!45, dotted, line width=0.5pt] (\FrenteX, 0) -- (\FrenteX, \CapaAltura);

\end{scope}
\end{tikzpicture}
\end{document}
